\section{Exact canonical residual reflection after projective cancellation}
\label{crr-section}

The general residual bill admits an exact formula when the initial
residual is power-free. Work in $\mathcal O=\mathbb Z[\zeta]$, where
$\zeta^2-\zeta+1=0$, and put $\gamma=1+\zeta$.
Fix an integer $n\ge2$ and
\[
 \tau=u\gamma^e v,\qquad e\in\{0,1\},\qquad V=N(v),
\]
where $u$ is a unit and $v$ is primitive and unramified. Let $w$ be
primitive and unramified with $R=N(w)>1$. In this section assume
that $V$ is $n$-free: $0\le v_p(V)<n$ for every rational prime $p$.
The case $V=1$ is included, and $V$ and $R$ need not be coprime.
Write
\[
 C=\gcd\bigl(|\operatorname{coord}_1(\tau w^n)|,
              |\operatorname{coord}_2(\tau w^n)|\bigr)>0,
 \qquad D=\operatorname{rad}(C).
\]
Primitive reduction gives the actual integer norm
\begin{equation}\label{crr-norm}
 M=N(\tau w^n/C)=3^eVR^n/C^2.
\end{equation}

\begin{theorem}[Exact canonical pair]\label{crr-pair}
There are unique positive integers $V_{\rm can},R_{\rm can}$ with
$V_{\rm can}$ $n$-free, $3\nmid V_{\rm can}R_{\rm can}$, and
$M=3^eV_{\rm can}R_{\rm can}^n$. They satisfy
\begin{gather}
 C\mid V,\qquad \gcd(C,V/C)=1,\qquad D\mid R,\label{crr-divisions}\\
 V_{\rm can}=\frac{VD^n}{C^2},\qquad
 R_{\rm can}=\frac RD.\label{crr-formulas}
\end{gather}
In particular, each canceled prime removes exactly one layer from
the canonical root norm, independently of its content depth.
\end{theorem}
\begin{proof}
Primitive unramified norm supports contain only split rational primes
$p\equiv1\pmod3$, hence $p\ge7$. The exponent of $\gamma$ in
$\tau w^n$ is exactly $e\le1$, so no rational factor $3$ is removed.
At a prime common to $V$ and $R$, set $a=v_p(V)$ and $f=v_p(R)$.
Each primitive factor chooses only one of the two prime elements
over $p$. If those choices agree, there is no rational $p$-content
and the reduced norm has depth $a+nf$.

If the choices are opposite, the rational content depth is
$\min(a,nf)$. Here $1\le a<n\le nf$, so this depth is exactly $a$,
and the reduced output depth is
\begin{equation}\label{crr-reflection}
 nf-a=(n-a)+n(f-1),\qquad 1\le n-a<n.
\end{equation}
Thus a canceled prime contributes its complete factor $p^a$ from
$V$ to $C$, contributes $p$ to $D$, changes the canonical residual
depth from $a$ to $n-a$, and changes the root depth from $f$ to $f-1$.
Every other prime keeps its residual and root depths, allowing zero
depths when it appears in only one input norm. This proves
\eqref{crr-divisions} and the primewise formulas
\eqref{crr-formulas}. In particular, $V/C$ and $D^n/C$ are coprime
integers whose positive prime exponents are less than $n$; their product
$VD^n/C^2$ is therefore an $n$-free integer. Together with
\eqref{crr-norm} this proves existence and uniqueness of the canonical
pair, with the first ramified contribution $3^e$ unchanged.
\end{proof}

The canonical root may be one. It is a nonunit precisely when $R>D$;
if $R=D$, every root prime has depth one and every root orientation
is canceled against its opposite residual orientation. The theorem
is a statement about the norm and the actual reduced output. It does
not assert that an oriented algebraic root is divided by the rational
integer $D$ inside $\mathcal O$.

\begin{corollary}[Exact residual bill]\label{crr-bill}
Under the same hypotheses,
\begin{equation}\label{crr-exact-bill}
 VV_{\rm can}=(V/C)^2D^n,
 \qquad
 \frac{\log V+\log V_{\rm can}}n
   =\log D+\frac2n\log(V/C).
\end{equation}
If $C>1$, this logarithmic bill is at least $\log D\ge\log7$.
Equality with $\log D$ holds exactly when $V=C$, and equality with
$\log7$ holds exactly when $D=7$ and $V=C$.
Every other positive integer decomposition
$M=3^eV'(R')^n$ has
\begin{equation}\label{crr-other-pair}
 V'=V_{\rm can}S^n,\qquad R_{\rm can}=SR'
\end{equation}
for a positive integer $S$.
\end{corollary}
\begin{proof}
The identity and its logarithm follow by substituting
\eqref{crr-formulas}. The equality criteria use $V/C\ge1$ and,
when $C>1$, the fact that all its primes are at least seven.
For the final assertion, each prime exponent in $V'$ has remainder
equal to its exponent in $V_{\rm can}$ modulo $n$. The quotient
is a nonnegative integer. Multiplying the corresponding prime powers
gives $S$;
the equality of norms then gives the root identity in
\eqref{crr-other-pair}. Thus a noncanonical residual can only
increase the canonical bill.
\end{proof}

The actual norm-seven family gives equality for every $n\ge2$.
With $\pi=2+\zeta$ and $w=\bar\pi=3-\zeta$, choose a unit $u_n$
so that $u_nw^{n-1}$ has positive coordinates. The identity
\[
 (u_n\pi)w^n=7u_nw^{n-1}
\]
has $V=R=C=D=7$, hence
$R_{\rm can}=1$ and $V_{\rm can}=7^{n-1}$.
The previously proved primitivity of every $w^m$ makes this the actual
content and the actual reduced norm. Its bill is exactly $\log7$.
It remains a single-map inverse example and does not assert a point
on a prescribed simultaneous double fiber product.

\begin{proposition}[The power-free premise cannot be omitted]
\label{crr-boundary}
The exact canonical-pair formula need not hold for a primitive
unramified multiplier whose norm is not $n$-free.
\end{proposition}
\begin{proof}
Keep $\pi=2+\zeta$ and let $n\ge3$. First take
$\tau=\pi^{n+1}$ and $w=\bar\pi$. Then
\[
 \tau w^n=7^n\pi,\qquad
 V=7^{n+1},\quad R=7,\quad C=7^n.
\]
Here $C$ no longer contains the complete prime-power factor of $V$,
so $\gcd(C,V/C)=1$ already fails, although the displayed canonical
pair formulas happen to give the correct pair $(7,1)$.

For an actual failure of the formulas themselves, take instead
$\tau=\pi^{n+1}$ and $w=\bar\pi^2$. Primitivity and the oriented
prime factorization give
\[
 V=7^{n+1},\qquad R=7^2,\qquad C=7^{n+1},\qquad
 N(\tau w^n/C)=7^{n-1}.
\]
The canonical pair is $(7^{n-1},1)$, whereas $R/D=7$ and
$VD^n/C^2=1/7$. All ring inputs remain primitive and unramified;
the failed hypothesis is precisely the $n$-freeness of $V$.
\end{proof}

This exact reflection result determines cancellation for supplied
actual inputs and their orientations. It does not construct rational
points on the simultaneous inverse curves, provide upper bounds for
their heights, or settle arbitrary-root signed-tail membership. The
general residual-divisibility bill continues to apply outside this
power-free subdomain. No Lean formalization of the exact canonical
pair formula is claimed here.
